\scanpage{196}

\[
\begin{aligned}
\underbrace{\varphi_{n'}\underbrace{\mathcal A_*(\tau)\mathcal A_*(\sigma)}_{\mathcal A_*(\sigma\tau)}}_{\mathcal A_*(\sigma\tau)\varphi\;(\text{p. hyp. sur }\varphi_{n'})}
&=\underbrace{\mathcal A_*(\tau)\underbrace{\varphi_n\mathcal A_*(\sigma)}_{\mathcal A_*(\sigma)\varphi\;(\text{par hyp. sur }\varphi_n)}}_{\mathcal A_*(\sigma\tau)\varphi}
\qquad\forall\sigma:\Delta_n\to\Delta_1
\end{aligned}
\]

OK.

\uline{Dém. de la prop} Le cas $n=0$ est trivial, pour le cas $n=1$ il faut considérer les trois cas $\sigma=\mathrm{id}$ qui donne $\varphi_1=\varphi$, puis $\sigma$ application constante de valeur $0$ ($\alpha_0$) ou $1$ ($\deleted{\unsure{\alpha}}\beta_0$), de sorte que $\mathcal A_*(\alpha_0)=\alpha$, $\mathcal A_*(\beta_0)=\beta$, et le fait que $\varphi_1=\varphi$ y soit comp. ce n'est autre que l'hyp. suppl. sur $\varphi$ sur les augmentations \inserted{$\alpha,\beta$}. On peut donc supposer $n\geq2$. Considérons les applications \inserted{\unsure{croiss|corresp.}} $\sigma:\Delta_n\to\Delta_1$ ; elles sont surj. ou constantes. Considérons les surjections, il y en a $n$, soit $\sigma_{\deleted{\unsure{\alpha}}r}$ ($0\leq\deleted{\unsure{\alpha}}r\leq n-1$)

\[
\sigma_{\deleted{\unsure{\alpha}}r}(0)\overset{\inserted{\sigma}}{=}0,\,-\,\sigma_{\deleted{\unsure{\alpha}}r}(\unsure{\alpha|r})=0,\quad
\sigma_{\deleted{\unsure{\alpha}}r}(\unsure{\alpha|r}+1)=1,\,-\,\sigma_r(n)=1
\]

donc \deleted{$\mathcal A$}

\[
\left\{\begin{aligned}
\mathcal A_*(\sigma_r)(X_0)&=\deleted{\unsure{U}}X_0+-+X_r\overset{\text{déf}}{=}U_r\\
\mathcal A_*(\sigma_r)(X_1)&=X_{r+1}\,-+X_n
\end{aligned}\right.
\]

On doit donc avoir

\[
\mathcal A_*(\sigma_r)(X_0^{(i)})=U_r^{(i)}
\]

\uline{Lemme} Les $U_r^{(i)}$ ($\deleted{\unsure{0\leq}}i\geq1,\ 0\leq r\leq n-1$) engendrent la \inserted{$S$-}algèbre $\mathcal A_n^0$, qui est isomorphe en fait à $\deleted{\unsure{k}}S\{U_0,-,U_{n-1}\}$,

% archive sheet number: 195
